\documentclass[11pt]{article}

\usepackage[margin=1in]{geometry}
\usepackage{amsmath,amssymb,amsthm}
\usepackage{booktabs}
\usepackage{array}
\usepackage[hidelinks]{hyperref}

\setlength{\emergencystretch}{1.5em}

\newtheorem{theorem}{Theorem}[section]
\newtheorem{proposition}[theorem]{Proposition}
\newtheorem{lemma}[theorem]{Lemma}
\newtheorem{corollary}[theorem]{Corollary}
\theoremstyle{definition}
\newtheorem{definition}[theorem]{Definition}
\newtheorem{remark}[theorem]{Remark}

\newcommand{\Cl}{\mathrm{Cl}}
\newcommand{\Mat}{M}
\newcommand{\F}{\mathbb{F}}
\newcommand{\Z}{\mathbb{Z}}
\newcommand{\Q}{\mathbb{Q}}
\newcommand{\R}{\mathbb{R}}
\newcommand{\C}{\mathbb{C}}
\newcommand{\HH}{\mathbb{H}}
\newcommand{\tr}{\mathrm{tr}}

\title{\textbf{Disproof and Exact Wedderburn--Artin Characterization of the\\
Minimal-Ideal-Count Idempotent Decomposition Conjecture}\\
\large (Resolution of Conjecture 00000007625)}
\author{Formalized in Core Lean 4.19.0 with Computational Verification}
\date{}

\begin{document}
\maketitle

\begin{abstract}
We resolve Conjecture~00000007625 in the negative across all standard
ring-theoretic and Clifford-algebraic interpretations, while proving the exact
Wedderburn--Artin central-block theorem that captures its valid core. The
conjecture asserts three linked principles for idempotent projection structures:
(I)~the count law of idempotent decompositions equals the minimal-ideal count;
(II)~minimality is realized by matrix factors; and (III)~the minimal realization
of matrix factors is the center decomposition. First, we prove a non-vacuous
positive theorem: in any unital ring $R$ whose identity decomposes into $t$
pairwise orthogonal simple central blocks $1 = c_1 + \dots + c_t$ (and in full
generality for split commutative semisimple algebras $K^t$, such as
$\Cl_{1,0}(K) \cong K \times K$ for $\mathrm{char}\,K \neq 2$), the unordered
central-primitive decomposition is unique, its summand length $t$ equals the
minimal two-sided ideal count, and on $K^t$ this also equals the minimal left
ideal count. Second, we show that as stated for general idempotent projection
structures, all three clauses fail:
(1)~on the split Clifford matrix algebra $\Cl_{1,1}(\F_3) \cong \Mat_2(\F_3)$,
the minimal two-sided ideal count ($1$), primitive decomposition length ($2$),
minimal left ideal count ($4$), unordered primitive decomposition count ($6$),
ordered primitive decomposition count ($12$), and total idempotent count ($14$)
are pairwise distinct, with $\{2,6,7,12,13,14\} \cap \{1,4,8\} = \emptyset$,
and over $\Z$ or $\Q$ there are infinitely many unordered primitive
decompositions while $m_{\mathrm{two}}(\Mat_2(\Q)) = 1$;
(2)~in division Clifford algebras $\Cl_{0,1}(\F_3) \cong \F_9$ and
$\Cl_{0,2}(K) = \HH(K)$ ($K \in \{\Q, \R\}$), the unique minimal ideal is a
division algebra non-isomorphic as a ring to $\Mat_n(K)$ for every $n \ge 1$;
and (3)~in the exterior algebra $\Lambda(\F_3^1) = \F_3[\varepsilon]/(\varepsilon^2)$
and the commutative local algebra $\F_3[x,y]/(x^2,xy,y^2)$, minimal ideals are
square-zero nilpotent (containing no nonzero idempotents) and even in the
commutative central setting satisfy $1 \neq 4$. All theorems are verified from
first principles in Core Lean~4.19.0 with zero external dependencies.
\end{abstract}

\section{Problem Statement and Taxonomic Analysis}

\subsection{Statement of the Conjecture}

Conjecture~00000007625 reads verbatim as follows:
\begin{quote}
\emph{Definition: Idempotents as projection structure. Conjecture: The count law
of idempotent decompositions is the minimal-ideal count; the realization of
minimality is the matrix factor; and the minimal realization of matrix factors
is the center decomposition. (minimal-ideal-count idempotent decomposition)}
\end{quote}
The accompanying Chinese formulation (\emph{idempotents as projection structure;
the number law of idempotent decompositions is the minimal-ideal count; the
realization of minimality is the matrix factor; and the minimal realization of
matrix factors is the center decomposition}) has the exact same three-clause
logical structure.

\subsection{Algebraic Definitions and Candidate Readings}

Let $R$ be a nonzero unital associative ring ($1 \neq 0$) with center
$Z(R) = \{z \in R : zx = xz \text{ for all } x \in R\}$.
\begin{itemize}
\item An element $e \in R$ is an \emph{idempotent} ($e \in E(R)$) if $e^2 = e$;
left multiplication $x \mapsto ex$ defines a module projection on every left
$R$-module and yields the Peirce decomposition $R = Re \oplus R(1-e)$.
\item Two elements $e, f \in R$ are \emph{orthogonal} if $ef = fe = 0$. A
nonzero idempotent $e \in R$ is \emph{primitive} if there do not exist nonzero
orthogonal idempotents $f, g \in R$ with $e = f + g$. A nonzero central
idempotent $c \in Z(R)$ is \emph{central-primitive} if there do not exist
nonzero orthogonal central idempotents $f, g \in Z(R)$ with $c = f + g$.
\item An \emph{orthogonal idempotent decomposition} of $1 \in R$ of
\emph{length} $r \ge 1$ is a tuple $(e_1, \dots, e_r)$ of nonzero pairwise
orthogonal idempotents satisfying $\sum_{i=1}^r e_i = 1$. It is \emph{primitive}
(complete) if every $e_i$ is primitive, and \emph{central-primitive} (the
\emph{center decomposition}) if every $e_i$ is central-primitive.
\item A nonzero left ideal $L \leq {}_R R$ (respectively, nonzero two-sided
ideal $I \trianglelefteq R$) is \emph{minimal} if every nonzero left ideal
$M \subseteq L$ (respectively, nonzero two-sided ideal $J \subseteq I$) equals
$L$ (respectively, $I$). We write $m_{\mathrm{left}}(R)$,
$m_{\mathrm{right}}(R)$, and $m_{\mathrm{two}}(R)$ for the cardinalities of the
sets of minimal nonzero left, right, and two-sided ideals of $R$.
\end{itemize}

Because both ``count law of idempotent decompositions'' and ``minimal-ideal
count'' admit several standard meanings in noncommutative ring theory, a
complete resolution must test all natural readings simultaneously:
\begin{itemize}
\item[(C1)] \textbf{Decomposition length $r(R)$:} the number $r$ of summands in a
primitive orthogonal decomposition $1 = e_1 + \dots + e_r$.
\item[(C2)] \textbf{Unordered decomposition count $N_{\mathrm{unord}}(R)$:} the number
of distinct sets $\{e_1, \dots, e_r\}$ arising from primitive orthogonal
decompositions of $1$ (or $N_{\mathrm{unord}}^{\mathrm{all}}(R)$ if
non-primitive orthogonal decompositions are also counted).
\item[(C3)] \textbf{Ordered decomposition count $N_{\mathrm{ord}}(R)$:} the number of
distinct ordered tuples $(e_1, \dots, e_r)$ of primitive orthogonal
decompositions of $1$ (or $N_{\mathrm{ord}}^{\mathrm{all}}(R)$ for all ordered
orthogonal decompositions).
\item[(C4)] \textbf{Idempotent element count $|E(R)|$:} the total number of
idempotents in $R$ (or $|E(R) \setminus \{0\}|$ for nonzero idempotents).
\end{itemize}
Likewise, ``minimal-ideal count'' may refer to $m_{\mathrm{two}}(R)$,
$m_{\mathrm{left}}(R)$ ($= m_{\mathrm{right}}(R)$ in self-opposite rings), or
even the cardinality of the union of minimal left and minimal right ideals.
Table~\ref{tab:invariants} summarizes all invariants across the benchmark
algebras analyzed in Sections~2--4.

\begin{table}[ht]
\centering
\footnotesize
\setlength{\tabcolsep}{4.5pt}
\caption{Exact invariants across semisimple, Clifford, division, and
non-semisimple benchmark algebras.}
\label{tab:invariants}
\vspace{4pt}
\begin{tabular}{@{}lccccccc@{}}
\toprule
Algebra $R$ & Center $t$ & $m_{\mathrm{two}}$ & $m_{\mathrm{left}}$ & Length $r$ & $N_{\mathrm{unord}}$ & $N_{\mathrm{ord}}$ & $|E(R)|$ \\
\midrule
$\F_3^t$ ($t \ge 1$) & $t$ & $t$ & $t$ & $t$ & $1$ & $t!$ & $2^t$ \\
$\Cl_{1,0}(\F_3) \cong \F_3 \times \F_3$ & $2$ & $2$ & $2$ & $2$ & $1$ & $2$ & $4$ \\
$\Mat_2(\F_2) \oplus \F_2$ & $2$ & $2$ & $4$ & $3$ & $3$ & $18$ & $16$ \\
$\Mat_2(\F_2)$ & $1$ & $1$ & $3$ & $2$ & $3$ & $6$ & $8$ \\
$\Cl_{1,1}(\F_3) \cong \Mat_2(\F_3)$ & $\mathbf{1}$ & $\mathbf{1}$ & $\mathbf{4}$ & $\mathbf{2}$ & $\mathbf{6}$ & $\mathbf{12}$ & $\mathbf{14}$ \\
$\Mat_2(\F_q)$ ($q$ prime power) & $1$ & $1$ & $q+1$ & $2$ & $\tfrac{1}{2}q(q+1)$ & $q(q+1)$ & $q(q+1)+2$ \\
$\Cl_{1,1}(\Q) \cong \Mat_2(\Q)$ & $1$ & $1$ & $\infty$ & $2$ & $\infty$ & $\infty$ & $\infty$ \\
$\Cl_{0,1}(\F_3) \cong \F_9$ ($\not\cong \Mat_n(\F_3)$) & $1$ & $1$ & $1$ & $1$ & $1$ & $1$ & $2$ \\
$\Cl_{0,2}(\Q) = \HH(\Q)$ ($\not\cong \Mat_n(\Q)$) & $1$ & $1$ & $1$ & $1$ & $1$ & $1$ & $2$ \\
$\Lambda(\F_3^1) = \F_3[\varepsilon]/(\varepsilon^2)$ ($I^2 = 0$) & $1$ & $1$ & $1$ & $1$ & $1$ & $1$ & $2$ \\
$\F_3[x,y]/(x^2,xy,y^2)$ (comm., $I^2 = 0$) & $1$ & $\mathbf{4}$ & $\mathbf{4}$ & $\mathbf{1}$ & $\mathbf{1}$ & $\mathbf{1}$ & $2$ \\
\bottomrule
\end{tabular}
\end{table}

\section{Positive Regime: Central Block Uniqueness and \texorpdfstring{$K^t$}{K\^t}}

We first establish the precise Wedderburn--Artin domain of validity behind
Conjecture~00000007625: when ``idempotent decomposition'' is restricted to
\emph{central-primitive} idempotents, ``count law'' means \emph{decomposition
length}, ``minimal ideal'' means \emph{minimal two-sided ideal}, and the ring
decomposes into simple central blocks $B_j \cong \Mat_{n_j}(D_j)$ over division
rings $D_j$.

\begin{definition}
Let $R$ be a unital ring. A central idempotent $c \in Z(R)$ is a \emph{simple
central block} if $c \neq 0$ and for every two-sided ideal $I \trianglelefteq R$,
whenever there exists $x \in I$ with $xc \neq 0$, one has $c \in I$. We denote
its associated block by $B(c) = \{x \in R : xc = x\} = Rc$.
\end{definition}

\begin{theorem}[Central Block Uniqueness and Minimal Two-Sided Ideal Count Law]
\label{thm:central-block}
Suppose $R$ admits a pairwise orthogonal decomposition $1 = c_1 + \dots + c_t$
into simple central blocks $c_1, \dots, c_t \in Z(R)$. Then:
\begin{enumerate}
\item[(a)] Each $B(c_j)$ is a minimal nonzero two-sided ideal of $R$, the blocks
$B(c_1), \dots, B(c_t)$ are pairwise distinct, and every minimal nonzero
two-sided ideal of $R$ equals $B(c_j)$ for a unique $j \in \{1,\dots,t\}$.
Consequently, $m_{\mathrm{two}}(R) = t$.
\item[(b)] For every central idempotent $z \in Z(R)$ and every $j \in \{1,\dots,t\}$,
either $zc_j = 0$ or $zc_j = c_j$. In particular, the set of central-primitive
idempotents of $R$ is precisely $\{c_1, \dots, c_t\}$, of cardinality $t$.
\item[(c)] Every central-primitive orthogonal decomposition $1 = d_1 + \dots + d_s$
satisfies $s = t = m_{\mathrm{two}}(R)$ and $\{d_1, \dots, d_s\} =
\{c_1, \dots, c_t\}$ (so the unordered center decomposition is unique).
\end{enumerate}
\end{theorem}

\begin{proof}
(a)~Because $c_j \in Z(R)$, $(xr)c_j = x(rc_j) = x(c_j r) = (xc_j)r = xr$ for
all $x \in B(c_j)$ and $r \in R$, so $B(c_j)$ is a two-sided ideal containing
$c_j \neq 0$. If $J \subseteq B(c_j)$ is a nonzero two-sided ideal containing
$x \neq 0$, then $xc_j = x \neq 0$, so the simple-block property yields
$c_j \in J$, hence $y = yc_j \in J$ for all $y \in B(c_j)$. Thus $B(c_j)$ is
minimal. For $i \neq j$, $c_i \in B(c_i)$ while $c_i c_j = 0 \neq c_i$, so
$B(c_i) \neq B(c_j)$. Conversely, if $I \trianglelefteq R$ is any minimal
nonzero two-sided ideal and $0 \neq x \in I$, then
$x = x \cdot 1 = \sum_{j=1}^t xc_j$ implies $xc_j \neq 0$ for some $j$, whence
$c_j \in I$ and $B(c_j) \subseteq I$, forcing $I = B(c_j)$ by minimality of $I$.

(b)~Let $z \in Z(R)$ satisfy $z^2 = z$ and put $w = zc_j$. Since $z$ and $c_j$
are commuting idempotents, $w$ is a central idempotent with $wc_j = w$, and
$B(w)$ is a two-sided ideal. If $w \neq 0$, then $w \in B(w)$ satisfies
$wc_j = w \neq 0$, so $c_j \in B(w)$, i.e.\ $c_j w = c_j$; on the other hand,
$c_j w = c_j z c_j = z c_j^2 = w$, giving $w = c_j$. Hence $zc_j \in \{0, c_j\}$.
If $c_j = f + g$ with nonzero orthogonal central idempotents $f, g$, then
$f = f(f+g) = fc_j \in \{0, c_j\}$ forces $f = c_j$ and $g = 0$, contradiction;
thus each $c_j$ is central-primitive. Conversely, if $d$ is central-primitive,
$d = \sum_{j=1}^t dc_j \neq 0$ yields $j$ with $dc_j \neq 0$, so $dc_j = c_j$.
Setting $g = d - c_j \in Z(R)$, one verifies $c_j g = g c_j = 0$ and $g^2 = g$;
since $d = c_j + g$ is central-primitive and $c_j \neq 0$, we have $g = 0$, so
$d = c_j$.

(c)~If $1 = d_1 + \dots + d_s$ is any central-primitive orthogonal decomposition,
part~(b) shows $\{d_1, \dots, d_s\} \subseteq \{c_1, \dots, c_t\}$. Conversely,
for each $j$, $c_j = \sum_{i=1}^s c_j d_i \neq 0$ gives some $i$ with
$c_j d_i \neq 0$; writing $d_i = c_k$ forces $j = k$ by orthogonality, so
$c_j = d_i$. Finally, pairwise orthogonality makes $i \mapsto d_i$ injective on
$\{1, \dots, s\}$, hence $s = |\{c_1, \dots, c_t\}| = t$.
\end{proof}

\begin{corollary}[Split Commutative Semisimple Algebras $K^t$ and $\Cl_{1,0}(K)$]
\label{cor:split-comm}
Let $K$ be a field and $R = K^t$ ($t \ge 1$) under componentwise operations,
with coordinate projections $e_j = (0,\dots,1,\dots,0)$. Then $R$ is
commutative, every $e_j$ is a simple central block with $B(e_j) \cong K \cong
\Mat_1(K)$, minimal left ideals coincide with minimal two-sided ideals, and
\[
  m_{\mathrm{two}}(K^t) = m_{\mathrm{left}}(K^t) = r(K^t) = t,
  \qquad N_{\mathrm{unord}}(K^t) = 1, \qquad N_{\mathrm{ord}}(K^t) = t!.
\]
In particular, when $\mathrm{char}\,K \neq 2$ (for example $K = \F_3$),
$\Cl_{1,0}(K) = K[e]/(e^2 - 1) \cong K \times K$ via $s + ve \mapsto (s+v, s-v)$
satisfies $m_{\mathrm{two}} = m_{\mathrm{left}} = r = 2$.
\end{corollary}

\begin{remark}
Even on $K^t$, Corollary~\ref{cor:split-comm} exposes an ambiguity in the phrase
``count law of idempotent decompositions'': while the \emph{length} $r(K^t) = t$
matches the minimal-ideal count $t$, the \emph{number of distinct decompositions}
does not ($N_{\mathrm{unord}}(K^t) = 1 \neq t$ for all $t \ge 2$, and
$N_{\mathrm{ord}}(K^3) = 3! = 6 \neq 3$). In Section~3 we show that once $R$ is
noncommutative, \emph{every} reading fails simultaneously.
\end{remark}

\section{Noncommutative Matrix and Split Clifford Refutation}

\subsection{Characteristic-Free Idempotent Theory on \texorpdfstring{$\Mat_2(D)$}{M\_2(D)}}

We begin with an elementary, characteristic-free characterization of idempotents
and primitive decompositions in $\Mat_2(D)$ valid over \emph{every} integral
domain $D$ (including $\mathrm{char}\,D = 2$).

\begin{proposition}[Idempotents and Primitive Decompositions in $\Mat_2(D)$]
\label{prop:mat2-domain}
Let $D$ be an integral domain and $M = \Mat_2(D)$.
\begin{enumerate}
\item[(a)] A matrix $m = \begin{pmatrix} a & b \\ c & d \end{pmatrix} \in M$ satisfies
$m^2 = m$ and $m \notin \{0, I_2\}$ if and only if $a + d = 1$ and $ad = bc$.
\item[(b)] If $f, g \in M$ are nonzero orthogonal idempotents ($f^2 = f \neq 0$,
$g^2 = g \neq 0$, $fg = gf = 0$), then $f + g = I_2$.
\item[(c)] An element $m \in M$ is a primitive idempotent if and only if
$a + d = 1$ and $ad = bc$; in particular, $I_2 = E_{11} + E_{22}$ is not
primitive.
\item[(d)] Every primitive orthogonal idempotent decomposition $I_2 = e_1 + \dots + e_r$
in $M$ has exact length $r = 2$, and $(e_1, e_2) \mapsto e_1$ is a bijection
from ordered primitive orthogonal decompositions of $I_2$ onto primitive
idempotents of $M$, with inverse $e \mapsto (e, I_2 - e)$.
\end{enumerate}
\end{proposition}

\begin{proof}
(a)~Suppose $a + d = 1$ and $ad = bc$. Direct expansion gives
$(m^2)_{11} = a^2 + bc = a(a+d) = a$, $(m^2)_{12} = b(a+d) = b$,
$(m^2)_{21} = c(a+d) = c$, and $(m^2)_{22} = cb + d^2 = (a+d)d = d$, so
$m^2 = m$. If $m = 0$ then $a + d = 0 \neq 1$, and if $m = I_2$ then
$1 = ad = bc = 0$, impossible since $1 \neq 0$. Conversely, suppose $m^2 = m$
and $m \notin \{0, I_2\}$. If $b \neq 0$, comparing $(1,2)$-entries gives
$b(a+d-1) = 0$, so $a + d = 1$ because $D$ is a domain, and $(m^2)_{11} = a$
then yields $a^2 + bc = a(a+d) = a^2 + ad$, hence $ad = bc$. The case $c \neq 0$
is identical via $(m^2)_{21} = c$. If $b = c = 0$, then $a^2 = a$ and $d^2 = d$
in the domain $D$ force $a, d \in \{0,1\}$; since $m \notin \{0, I_2\}$,
$(a,d) \in \{(1,0),(0,1)\}$, both satisfying $a + d = 1$ and $ad = 0 = bc$.

(b)~Neither $f$ nor $g$ can equal $I_2$ (for if $f = I_2$, then $g = I_2 g = fg = 0$).
By~(a), $f_{11} + f_{22} = 1$ and $g_{11} + g_{22} = 1$. Set $h = f + g$. Since
$fg = gf = 0$, we have $h^2 = f^2 + fg + gf + g^2 = h$. If $h = 0$, then
$f = f^2 + fg = fh = 0$, contradiction. If $h \neq I_2$, then~(a) applied to $h$
gives $h_{11} + h_{22} = 1$; yet $h_{11} + h_{22} = (f_{11} + f_{22}) +
(g_{11} + g_{22}) = 1 + 1$, whence $1 + 1 = 1$ and $1 = 0$ in $D$, contradiction.
Therefore $f + g = h = I_2$.

(c)~Clearly $I_2 = E_{11} + E_{22}$ splits into nonzero orthogonal idempotents,
so $I_2$ is not primitive; thus every primitive idempotent satisfies
$m \notin \{0, I_2\}$ and hence $a + d = 1$, $ad = bc$ by~(a). Conversely, if
$a + d = 1$ and $ad = bc$, then $m \neq I_2$, and any splitting $m = f + g$ into
nonzero orthogonal idempotents would force $m = f + g = I_2$ by~(b), contradiction.

(d)~Let $I_2 = e_1 + \dots + e_r$ be a primitive orthogonal decomposition. We
cannot have $r = 0$ ($0 \neq I_2$) or $r = 1$ ($e_1 = I_2$ is not primitive by~(c)).
If $r \ge 3$, then $e_1 + e_2 = I_2$ by~(b), and multiplying on the right by
$e_3$ yields $e_3 = I_2 e_3 = (e_1 + e_2)e_3 = e_1 e_3 + e_2 e_3 = 0$,
contradicting $e_3 \neq 0$. Hence $r = 2$, $e_2 = I_2 - e_1$, and conversely for
every primitive idempotent $e$, $I_2 - e$ has trace $(1-a) + (1-d) = 1$ and
determinant $(1-a)(1-d) - bc = 1 - (a+d) + (ad - bc) = 0$ with
$e(I_2 - e) = (I_2 - e)e = 0$.
\end{proof}

\subsection{Master Six-Way Separation on \texorpdfstring{$\Cl_{1,1}(\F_3) \cong \Mat_2(\F_3)$}{Cl\_\{1,1\}(F\_3) = M\_2(F\_3)}}

\begin{proposition}[Simplicity, Center, and Finite-Field Enumeration]
\label{prop:mat2-field}
Let $K$ be any field.
\begin{enumerate}
\item[(a)] $\Mat_2(K)$ is a simple ring ($m_{\mathrm{two}}(\Mat_2(K)) = 1$), its
center is $Z(\Mat_2(K)) = K \cdot I_2$, and $I_2$ is its unique central-primitive
idempotent ($t = 1 \neq 2 = r$).
\item[(b)] When $K = \F_q$ is a finite field of order $q$, the invariants of
$\Mat_2(\F_q)$ are:
\[
\begin{aligned}
  m_{\mathrm{two}} &= 1, \qquad r = 2, \qquad m_{\mathrm{left}} = m_{\mathrm{right}} = q+1,\\
  N_{\mathrm{unord}} &= \tfrac{1}{2}q(q+1), \qquad N_{\mathrm{ord}} = q(q+1), \qquad
  |E| = q(q+1)+2.
\end{aligned}
\]
\end{enumerate}
\end{proposition}

\begin{proof}
(a)~Let $I \trianglelefteq \Mat_2(K)$ contain $m \neq 0$. If $m_{uv} \neq 0$ for
$u, v \in \{1,2\}$, then $E_{1u} m E_{v1} + E_{2u} m E_{v2} = m_{uv} I_2 \in I$;
multiplying by $m_{uv}^{-1} I_2$ gives $I_2 \in I$, hence $I = \Mat_2(K)$.
Commuting with $E_{11}$ and $E_{12}$ forces any $z \in Z(\Mat_2(K))$ to satisfy
$z_{12} = z_{21} = 0$ and $z_{11} = z_{22}$, so $Z(\Mat_2(K)) \cong K$ has only
$0$ and $I_2$ as central idempotents.

(b)~By Proposition~\ref{prop:mat2-domain}(c), primitive idempotents are matrices
with $d = 1 - a$ and $bc = a(1-a)$:
\begin{itemize}
\item If $a \in \{0, 1\}$ ($2$ choices), $bc = 0$ has $2q - 1$ solutions
$(b,c) \in \F_q^2$, contributing $4q - 2$ matrices.
\item If $a \in \F_q \setminus \{0, 1\}$ ($q - 2$ choices), $a(1-a) \neq 0$, so
$b \in \F_q^\times$ ($q - 1$ choices) uniquely determines $c = a(1-a)b^{-1}$,
contributing $(q-2)(q-1)$ matrices.
\end{itemize}
Summing gives $(4q-2) + (q^2 - 3q + 2) = q(q+1)$ primitive idempotents (hence
$N_{\mathrm{ord}} = q(q+1)$ and $|E| = q(q+1) + 2$ after adjoining $0$ and $I_2$).
Since $e = I_2 - e$ would imply $0 = e(I_2 - e) = e^2 = e$, impossible, the
involution $e \mapsto I_2 - e$ is fixed-point-free and partitions primitive
idempotents into $\frac{1}{2}q(q+1)$ unordered decompositions $\{e, I_2 - e\}$.
Finally, minimal left ideals of $\Mat_2(\F_q)$ correspond bijectively to the
$q + 1$ points $[u : v] \in \mathbb{P}^1(\F_q)$ via
$L_{[u:v]} = \{m \in \Mat_2(\F_q) : \text{each row of } m \text{ lies in } \F_q(u,v)\}$.
\end{proof}

Notice that at $q = 2$, $\frac{1}{2}q(q+1) = 3$ happens by sheer numerical
coincidence to equal $q + 1 = 3$ (though $r = 2 \neq 3$ and $m_{\mathrm{two}} = 1$).
Setting $q = 3$ eliminates even this accidental coincidence and embeds directly
into Clifford algebras since $\mathrm{char}\,\F_3 = 3 \neq 2$.

\begin{theorem}[Master Six-Way Separation on $\Cl_{1,1}(\F_3) \cong \Mat_2(\F_3)$]
\label{thm:master-f3}
Let $\Cl_{1,1}(\F_3)$ be the Clifford algebra over $\F_3^2$ generated by
$\gamma_1, \gamma_2$ subject to $\gamma_1^2 = 1$, $\gamma_2^2 = -1$, and
$\gamma_1 \gamma_2 + \gamma_2 \gamma_1 = 0$.
\begin{enumerate}
\item[(a)] The map $s + v_1 \gamma_1 + v_2 \gamma_2 + v_{12}\gamma_1\gamma_2 \mapsto
\begin{pmatrix} s + v_1 & v_2 + v_{12} \\ -v_2 + v_{12} & s - v_1 \end{pmatrix}$
is an explicit isomorphism of unital rings $\Cl_{1,1}(\F_3) \cong \Mat_2(\F_3)$.
\item[(b)] On $R = \Mat_2(\F_3) \cong \Cl_{1,1}(\F_3)$, the algebraic invariants take
the strictly increasing values
\[
  m_{\mathrm{two}}(R) = 1 < r(R) = 2 < m_{\mathrm{left}}(R) = 4
  < N_{\mathrm{unord}}(R) = 6 < N_{\mathrm{ord}}(R) = 12 < |E(R)| = 14.
\]
Moreover, including non-primitive decompositions gives
$N_{\mathrm{unord}}^{\mathrm{all}}(R) = 7$,
$N_{\mathrm{ord}}^{\mathrm{all}}(R) = 13$, and $|E(R) \setminus \{0\}| = 13$,
while $m_{\mathrm{right}}(R) = 4$ is disjoint from $m_{\mathrm{left}}(R)$ so
$|\mathcal{I}_{\mathrm{left}}(R) \cup \mathcal{I}_{\mathrm{right}}(R)| = 8$.
Because
\[
  \{2, 6, 7, 12, 13, 14\} \cap \{1, 4, 8\} = \emptyset,
\]
\textbf{Clause~I of Conjecture~00000007625 is false under every interpretation
of both the decomposition count law and the minimal-ideal count.}
\end{enumerate}
\end{theorem}

\begin{proof}
Part~(a) follows by direct verification that $\gamma_1 \mapsto \mathrm{diag}(1,2)$
and $\gamma_2 \mapsto \begin{pmatrix} 0 & 1 \\ 2 & 0 \end{pmatrix}$ satisfy the
defining Clifford relations over $\F_3$ and have explicit two-sided inverse
$\begin{pmatrix} a & b \\ c & d \end{pmatrix} \mapsto
2(a+d) + 2(a-d)\gamma_1 + 2(b-c)\gamma_2 + 2(b+c)\gamma_1\gamma_2$ (using
$2^{-1} = 2$ in $\F_3$). Part~(b) specializes Proposition~\ref{prop:mat2-field}
to $q = 3$: the four minimal left ideals are explicitly $L_{[1:0]}$, $L_{[0:1]}$,
$L_{[1:1]}$, and $L_{[1:2]}$, none of which is a right ideal (since right
multiplication by $\Mat_2(\F_3)$ sweeps out all of $\Mat_2(\F_3)$), and the
$12$ primitive idempotents pair into $6$ complementary sets $\{e, I_2 - e\}$.
\end{proof}

\begin{remark}[Structural Failure of Clauses~II and~III on $\Mat_2(K)$]
Even on $\Mat_2(K)$, Clauses~II and~III suffer from a fundamental conflation
between one-sided projections and two-sided blocks: a primitive idempotent
$e_i$ generates a minimal \emph{left} ideal $\Mat_2(K)e_i$ (which contains
nonzero nilpotent elements such as $E_{21} \in \Mat_2(K)E_{11}$ with
$E_{21}^2 = 0$, hence is not even a semiprime ring, let alone a matrix ring),
whereas the matrix factor $\Mat_2(K)$ itself is cut out by the central
idempotent $I_2$ of length $1 \neq 2$. Similarly, in the two-block ring
$R = \Mat_2(\F_2) \oplus \F_2$, the center decomposition has length
$t = m_{\mathrm{two}}(R) = 2$, whereas primitive idempotent decompositions have
length $r = 3$, $N_{\mathrm{unord}} = 3$, $N_{\mathrm{ord}} = 18$, and
$m_{\mathrm{left}} = 4$.
\end{remark}

\begin{proposition}[Infinite Decompositions over $\Z$ and Characteristic-Zero Fields]
\label{prop:mat2-inf}
In $\Mat_2(\Z)$ (and in $\Mat_2(K) \cong \Cl_{1,1}(K)$ for any characteristic-zero
field $K$, such as $\Q$ or $\R$), for every $t \in \Z$ the matrices
\[
  e_t = \begin{pmatrix} 1 & t \\ 0 & 0 \end{pmatrix}, \qquad
  f_t = I_2 - e_t = \begin{pmatrix} 0 & -t \\ 0 & 1 \end{pmatrix}
\]
form a primitive orthogonal idempotent decomposition of $I_2$, and
$\{e_t, f_t\} \neq \{e_s, f_s\}$ whenever $t \neq s$. Thus $\Mat_2(\Z)$ and
$\Mat_2(\Q)$ admit infinitely many unordered primitive decompositions, whereas
$m_{\mathrm{two}}(\Mat_2(\Q)) = 1$ and $r(\Mat_2(\Q)) = 2$.
\end{proposition}

\begin{proof}
Both $e_t$ and $f_t$ have trace $1$ and determinant $0$ over the domain $\Z$
(or $K$), so they are primitive idempotents by Proposition~\ref{prop:mat2-domain}(c),
with $e_t + f_t = I_2$ and $e_t f_t = f_t e_t = 0$. If $\{e_t, f_t\} = \{e_s, f_s\}$,
comparing $(1,1)$-entries ($1 \neq 0$) forces $e_t = e_s$, hence $t = s$.
\end{proof}

\section{Division-Algebra and Non-Semisimple Obstructions}

We now show that Clause~II (``the realization of minimality is the matrix
factor'') and Clause~III (``the minimal realization of matrix factors is the
center decomposition'') also fail independently of noncommutativity of matrix
units: first via non-split division Clifford algebras, and second via degenerate
Clifford and commutative non-semisimple algebras.

\begin{theorem}[Division-Algebra Obstructions: $\Cl_{0,1}(\F_3)$ and $\Cl_{0,2}(K)$]
\label{thm:division-obstruction}
\begin{enumerate}
\item[(a)] The Clifford algebra $\Cl_{0,1}(\F_3) = \F_3[i]/(i^2 + 1) \cong \F_9$ is a
simple ring with $m_{\mathrm{two}} = m_{\mathrm{left}} = 1$ and $|E| = 2$, yet
for every integer $n \ge 1$ there is \textbf{no ring isomorphism}
$\Cl_{0,1}(\F_3) \cong \Mat_n(\F_3)$.
\item[(b)] Let $C$ be a commutative ring and $\HH(C) = \Cl_{0,2}(C)$ the quaternion
algebra with basis $1, \mathbf{i}, \mathbf{j}, \mathbf{k}$ satisfying
$\mathbf{i}^2 = \mathbf{j}^2 = -1$ and $\mathbf{i}\mathbf{j} = \mathbf{k} =
-\mathbf{j}\mathbf{i}$. For $C = \Z$, there is \textbf{no ring isomorphism}
$\HH(\Z) \cong \Mat_n(\Z)$ for any $n \ge 1$. Moreover, over any formally real
field $K$ (meaning $1 + 1 \neq 0$ and $a^2 + b^2 + c^2 + d^2 = 0 \Rightarrow
a = b = c = d = 0$, such as $K = \Q$ or $K = \R$), $\HH(K) = \Cl_{0,2}(K)$ is a
simple division ring ($m_{\mathrm{two}} = m_{\mathrm{left}} = 1$, $|E| = 2$)
that is \textbf{not ring-isomorphic to $\Mat_n(K)$ for any $n \ge 1$}.
\end{enumerate}
\end{theorem}

\begin{proof}
(a)~In $\F_3$, $a^2 + b^2 = 0$ implies $a = b = 0$, so every nonzero $a + bi \in
\Cl_{0,1}(\F_3)$ has inverse $(a - bi)(a^2 + b^2)^{-1}$; hence $\Cl_{0,1}(\F_3)$
is a field of order $9$. For $n = 1$, $|\Mat_1(\F_3)| = 3 \neq 9$, so no
bijection exists. For $n \ge 2$, $\Mat_n(\F_3)$ contains the nonzero nilpotent
matrix $E_{12}$ ($E_{12}^2 = 0$), whereas a field has no nonzero zero divisors.

(b)~For any $q = a + b\mathbf{i} + c\mathbf{j} + d\mathbf{k} \in \HH(C)$ and
$\bar{q} = a - b\mathbf{i} - c\mathbf{j} - d\mathbf{k}$, direct Cayley--Dickson
computation gives $\bar{q}q = q\bar{q} = (a^2 + b^2 + c^2 + d^2) \cdot 1$.
Over $\Z$ (respectively, a formally real field $K$), $q \neq 0$ implies
$N(q) = a^2 + b^2 + c^2 + d^2 \neq 0$, so $\HH(\Z)$ has no nonzero zero divisors
and $\HH(K)$ is a division ring (with $q^{-1} = N(q)^{-1}\bar{q}$, hence simple).
For $n = 1$, $\Mat_1(\Z) \cong \Z$ and $\Mat_1(K) \cong K$ are commutative,
whereas $\mathbf{i}\mathbf{j} = \mathbf{k} \neq -\mathbf{k} = \mathbf{j}\mathbf{i}$
in $\HH(\Z)$ and $\HH(K)$ (since $1 + 1 \neq 0$). For $n \ge 2$, $\Mat_n(\Z)$
and $\Mat_n(K)$ contain the nonzero nilpotent $E_{12}$ ($E_{12}^2 = 0$),
precluding any ring isomorphism with a zero-divisor-free ring.
\end{proof}

Thus even for finite-dimensional semisimple $K$-algebras, minimal two-sided
ideals are Wedderburn blocks $\Mat_n(D)$ over \emph{division $K$-algebras} $D$,
not split matrix rings $\Mat_n(K)$ over the ground field $K$. Finally, when
semisimplicity (or nondegeneracy of the quadratic form $Q$ in $\Cl(V,Q)$) is
dropped, all three clauses collapse even in the commutative case.

\begin{theorem}[Degenerate Clifford and Commutative Non-Semisimple Counterexamples]
\label{thm:non-semisimple}
The following two finite $\F_3$-algebras refute Clauses~I--III outside the
semisimple regime:
\begin{enumerate}
\item[(a)] \textbf{Exterior algebra.} Let $A = \Lambda(\F_3^1) =
\F_3[\varepsilon]/(\varepsilon^2)$ be the $Q = 0$ degenerate Clifford algebra.
Its only idempotents are $0$ and $1$, and its unique minimal nonzero two-sided
ideal is $I = \F_3 \varepsilon$ ($m_{\mathrm{two}} = 1$). However, $I^2 = 0$, so
$I$ contains \emph{no} nonzero idempotent, and $A$ is not ring-isomorphic to
$\Mat_n(\F_3)$ for any $n \ge 1$ (refuting Clauses~II and~III).
\item[(b)] \textbf{Commutative local algebra.} Let $R = \F_3[x,y]/(x^2, xy, y^2)$.
Then $R$ is commutative ($Z(R) = R$, and left ideals coincide with two-sided
ideals) with $E(R) = \{0, 1\}$, so $1$ is primitive and central-primitive and
\[
  r(R) = N_{\mathrm{unord}}(R) = N_{\mathrm{ord}}(R) = 1.
\]
Yet $R$ has \textbf{four} minimal nonzero two-sided ideals---namely $\F_3 x$,
$\F_3 y$, $\F_3(x+y)$, and $\F_3(x+2y)$---each satisfying $I^2 = 0$, so
\[
  m_{\mathrm{two}}(R) = m_{\mathrm{left}}(R) = 4 \neq 1.
\]
\end{enumerate}
\end{theorem}

\begin{proof}
(a)~If $(a + b\varepsilon)^2 = a^2 + 2ab\varepsilon = a + b\varepsilon$ in
$\F_3[\varepsilon]/(\varepsilon^2)$, then $a \in \{0,1\}$; both $a = 0$ ($0 = b$)
and $a = 1$ ($2b = b \Rightarrow b = 0$) force $b = 0$. Every element with
$a \neq 0$ has inverse $a^{-1} - a^{-2}b\varepsilon$, so any proper nonzero ideal
lies in $I = \F_3\varepsilon$, which is $1$-dimensional over $\F_3$ and hence the
unique minimal nonzero ideal, with $(b\varepsilon)(b'\varepsilon) = 0$. Finally,
$|\F_3[\varepsilon]/(\varepsilon^2)| = 9 \neq 3 = |\Mat_1(\F_3)|$, while
$\Mat_n(\F_3)$ is noncommutative for $n \ge 2$ ($E_{11}E_{12} = E_{12} \neq 0 =
E_{12}E_{11}$).

(b)~In $R = \F_3[x,y]/(x^2,xy,y^2)$, $(a + bx + cy)^2 = a^2 + 2abx + 2acy$ equals
$a + bx + cy$ iff $a \in \{0,1\}$ and $b = c = 0$, so $E(R) = \{0,1\}$ and all
decomposition lengths and counts equal $1$. Every element with $a \neq 0$ has
unit inverse $a^{-1} - a^{-2}bx - a^{-2}cy$, so every proper nonzero ideal lies
in the square-zero Jacobson radical $J = \F_3 x \oplus \F_3 y$. Because
$(a + bx + cy)(u x + v y) = a(u x + v y)$ acts on $J$ solely by $\F_3$-scalar
multiplication, the minimal nonzero ideals of $R$ are precisely the
$3 + 1 = 4$ one-dimensional $\F_3$-subspaces of $J$.
\end{proof}

\section{Formal Verification in Core Lean 4.19.0 and Computational Suite}

The entire mathematical development is formalized from first principles in Core
Lean~4.19.0 (\texttt{package/lean/}, zero external dependencies, compiled with
\texttt{-DwarningAsError=true}, zero \texttt{sorry}, zero \texttt{native\_decide},
and zero custom axioms). Exact finite cardinalities of predicates $P : \alpha \to
\mathrm{Prop}$ (including ideal predicates $I : R \to \mathrm{Prop}$) are
formalized uniformly via \texttt{ExactCard P k}, with constructive Pigeonhole
uniqueness (\texttt{exactCard\_unique : ExactCard P m $\to$ ExactCard P n $\to$
m = n}). Table~\ref{tab:lean-map} records the correspondence between paper
results and Lean~4 declarations, all assembled in the top-level theorem
\texttt{conjecture\_00000007625\_resolution}.

\begin{table}[ht]
\centering
\footnotesize
\setlength{\tabcolsep}{5pt}
\caption{Correspondence between mathematical theorems and Core Lean~4.19.0 declarations.}
\label{tab:lean-map}
\vspace{4pt}
\begin{tabular}{@{}lll@{}}
\toprule
Paper Result & Module & Lean~4 Declaration \\
\midrule
Pigeonhole uniqueness of \texttt{ExactCard} & \texttt{Basic} & \texttt{exactCard\_unique} \\
Theorem~\ref{thm:central-block}(a)--(c) & \texttt{Positive} & \texttt{central\_block\_minimal\_ideals\_exactCard} \\
 & \texttt{Positive} & \texttt{central\_block\_decomp\_uniqueness} \\
Corollary~\ref{cor:split-comm} ($\F_3^{m+1}$ \& $\Cl_{1,0}(\F_3)$) & \texttt{Positive} & \texttt{f3Pow\_positive\_theorem} \\
 & \texttt{Positive} & \texttt{cl10\_f3\_iso\_f3Pow2} \\
Proposition~\ref{prop:mat2-domain}(a)--(d) & \texttt{MatrixCounterexample} & \texttt{mat2\_prim\_idemp\_iff} \\
 & \texttt{MatrixCounterexample} & \texttt{mat2\_prim\_decomp\_length\_eq\_two} \\
Proposition~\ref{prop:mat2-field}(a) & \texttt{MatrixCounterexample} & \texttt{mat2\_isSimpleRing} \\
 & \texttt{MatrixCounterexample} & \texttt{mat2\_central\_prim\_exactCard\_one} \\
Theorem~\ref{thm:master-f3}(a)--(b) & \texttt{MatrixCounterexample} & \texttt{cl11\_f3\_iso\_mat2\_f3} \\
 & \texttt{MatrixCounterexample} & \texttt{mat2\_f3\_master\_refutation} \\
Proposition~\ref{prop:mat2-inf} ($\Mat_2(\Z)$) & \texttt{MatrixCounterexample} & \texttt{mat2\_int\_infinite\_unord\_prim\_decomps} \\
Theorem~\ref{thm:division-obstruction}(a) ($\Cl_{0,1}(\F_3)$) & \texttt{Obstructions} & \texttt{cl01\_f3\_not\_iso\_any\_matN\_f3} \\
Theorem~\ref{thm:division-obstruction}(b) ($\HH(\Z)$ \& $\HH(K)$) & \texttt{Obstructions} & \texttt{quatInt\_not\_iso\_any\_matN\_int} \\
 & \texttt{Obstructions} & \texttt{quatField\_simple\_and\_not\_matN} \\
Theorem~\ref{thm:non-semisimple}(a) ($\F_3[\varepsilon]/(\varepsilon^2)$) & \texttt{Obstructions} & \texttt{dualF3\_obstruction\_theorem} \\
Theorem~\ref{thm:non-semisimple}(b) ($\F_3[x,y]/(x^2,xy,y^2)$) & \texttt{Obstructions} & \texttt{localF3\_commutative\_counterexample} \\
End-to-end resolution bundle & \texttt{Conjecture00000007625} & \texttt{conjecture\_00000007625\_resolution} \\
\bottomrule
\end{tabular}
\end{table}

In addition, the standalone Python~3.12 suite \texttt{package/verify.py}
independently checks ideal closures, center elements, idempotent counts,
orthogonal decompositions, Clifford relations, quaternion inverses, and rational
infinite families across $\F_3^t$, $\Cl_{1,0}(\F_3)$, $\Mat_2(\F_2) \oplus \F_2$,
$\Mat_2(\F_q)$ ($q \in \{2,3,5,7\}$), $\Cl_{1,1}(\Q)$, $\Cl_{0,1}(\F_3)$,
$\Cl_{0,2}(\Q)$, $\F_3[\varepsilon]/(\varepsilon^2)$, $\F_3[x,y]/(x^2,xy,y^2)$,
and $T_2(\F_3)$.

\begin{thebibliography}{9}

\bibitem{Lam2001}
T.~Y.~Lam, \emph{A First Course in Noncommutative Rings}, 2nd~ed., Graduate
Texts in Mathematics, vol.~131, Springer, New York, 2001.

\bibitem{Lam2005}
T.~Y.~Lam, \emph{Introduction to Quadratic Forms over Fields}, Graduate Studies
in Mathematics, vol.~67, American Mathematical Society, Providence, RI, 2005.

\bibitem{LawsonMichelsohn1989}
H.~B.~Lawson,~Jr.\ and M.-L.~Michelsohn, \emph{Spin Geometry}, Princeton
Mathematical Series, vol.~38, Princeton University Press, Princeton, NJ, 1989.

\end{thebibliography}

\end{document}
